\documentclass[12pt]{article}
\newtheorem{theorem}{Proposition}
\usepackage {amsfonts}

\begin{document}

\title{\large Ces\`{a}ro Sum Approximation of \\ Outer Functions}

\author{\large R.W. Barnard, K. Pearce \\ \normalsize Department of Mathematics and Statistics \\ 
\normalsize Texas Tech University \\
\normalsize barnard@math.ttu.edu, pearce@math.ttu.edu
\and \large J. Cima \\ \normalsize Department of Mathematics \\
\normalsize University of North Carolina at Chapel Hill \\
\normalsize cima@math.unc.edu }

\date{}
\maketitle 

\begin{abstract}
 It is well known that outer functions are zero-free on the unit disk.  If an outer function, $f$, is given as an 
infinite series and a finite (polynomial) approximation is chosen, then it is desirable that  the 
approximants retain the zero-free property of $f$.   We observe for outer functions that the standard Taylor 
approximants do not, in general, retain the zero-free property -- even 
when fairly restrictive conditions are placed on the permissible outer functions.    We show, using methods 
of geometric function theory, that Ces\`{a}ro sum approximants for 
outer functions which arise as the derivatives of bounded convex functions do inherit the desired zero-free property .  
We, also, find that a ``cone-like" condition holds for the boundaries of the ranges of these approximants.

\end{abstract}



\section{\large Introduction}

Let $\Bbb D$ denote the open unit disk in $\Bbb C$.    It is well known that outer functions are zero-free on the unit disk.  
Outer functions, which play an important role in $H^p$ theory, arise in the characteristic equation which 
determines the stability of certain nonlinear systems of differential equations.   The solutions of such 
characteristic equations frequently involves ratios of the form $h/f$ where $f$ is an outer function 
(see \cite{Cu}, p.  288).  If $f$ is given as an infinite series and a finite (polynomial) approximation is 
chosen, then it is desirable, in order to justify the choice of the approximant, that the zeros of the 
approximant retain  the zero-free property of $f$.  

In this note we consider questions about when approximating sequences of polynomials for outer 
functions inherit the zero-free property on $\Bbb D$ possessed by outer functions.  This leads to an 
investigation of the location of the zeros and behavior of such approximating sequences of 
polynomials.

We show by examples that for outer functions the standard Taylor approximants do not, in general, 
retain the zero-free property.  This is especially the case for low order Taylor approximants, even 
when fairly restrictive conditions are placed on the permissible outer functions, such as requiring 
these to be generated by Smirnov domains or requiring these to satisfy certain geometric conditions.  
The case for high order Taylor approximants to functions analytic and non-vanishing on the closed 
unit disk is covered by Hurwitz's theorem which assures that for $n$, the degree of the approximant, 
sufficiently large, that the Taylor approximant will be zero-free on $\Bbb D$ since the target, the outer function, 
is zero-free on $\Bbb D$.  For practical purposes, however, the $n$ required above may be prohibitively large.  
Thus, we seek conditions on outer functions and classes of approximants for which the zero-free 
property on $\Bbb D$ will be inherited by all of the approximants.  

	Jentzsch's classical results (\cite{Di}, p. 352) show that the circle of convergence for a Taylor series 
is a subset of the set of limit points of the zero sets of the sequence of Taylor approximants.  This 
suggests that, in general, a careful analysis will be required in order to affirm the desired zero-free 
inheritance.

	We, thus, consider other approximants such as Ces\`{a}ro sums.  We find that when considering 
geometric restrictions, such as convexity on Smirnov domains, the methods of geometric function 
theory can be applied to verify the desired zero-free inheritance for appropriate approximants.  We, 
also, find that a ``cone-like" condition holds for the boundaries of the ranges of these approximants.

\section{\large Outer Functions}


	Recall that an outer function is a function $f$ in $H^p$ of the form

\begin{equation}
f(z) = e^{i \gamma} e^{ \frac{1} {2 \pi} \int _{-\pi} ^ {\pi}
\frac{1 + e^{it}z} {1-e^{it}z} \log \psi (t)dt }
\end{equation}


\noindent where  $\psi(t) \geq 0$,    $\log  \psi(t)$   is in $L^1$ and  $\psi(t)$ is in $L^p$.  See \cite{Du1} for the definitions and 
classical properties of outer functions.  Since any function $f$ in $H^1$ which has $1/f$ in $H^1$ is an outer function, 
then typical examples of outer functions can be generated by functions of the form 
$\prod _ {k=1} ^ n (1 - e^{i \theta _ k}z)^{\alpha _ k}$ for $-1 <  \alpha _k < 1$.  


\section{\large Taylor Approximants}

	We will look at some examples which show that the standard Taylor approximants are not necessarily 
zero-free on $\Bbb D$.  Consider $f(z) = (1+e^{i/10}z)^{1/2} /(1-z)^{71/100}$.  If we let 
$p_n(z) = \sum _{k=0} ^ n a_k z^k$  be the Taylor approximant for 
$f(z) = \sum _{k=0} ^ \infty  a_k z^k$, then computer computations show that $p_3$ has $1$ zero inside $\Bbb D$, 
$p_4$ has no zeros inside $\Bbb D$ and $p_5$ again has a zero inside $\Bbb D$.  Thus, not only do some of the $p_n$'s 
have zeros inside $\Bbb D$, there is also no guarantee that once they are zero-free on $\Bbb D$, that they will remain so for 
higher orders.

	Next, we show that even imposing restrictive geometric constraints on the outer functions may not suffice to 
assure that their Taylor approximants will inherit their zero-free property on $\Bbb D$.  A common and useful generation for 
outer functions is found by considering Smirnov domains.  See \cite{Po} for their definition and properties.  We recall that
 if a simply connected domain $G$ with a rectifiable boundary is a Smirnov domain, then the conformal map $f$, 
mapping $\Bbb D$ onto $G$, has the property that its derivative $f'$  is an outer function.  Also, boundedness and 
convexity is a sufficient condition on a domain to guarantee it to be a Smirnov domain.  Thus, the derivatives of 
bounded convex functions are always outer functions.

	For $r, \ 0 < r < 1$, let $f_r(z) = z/(1-rz)$.  Then,  $f_r$ maps $\Bbb D$ onto a bounded convex domain so that $f_r$  
is an outer function.  Since $f_r (z) = 1 + 2rz + ... $ it is easy to see that many of its early Taylor approximants are 
not zero-free on $\Bbb D$ for $r$ near $1$.

\section{\large Ces\`{a}ro Approximants}

	The previous observations suggest using different approximants.  We consider the Ces\`{a}ro means for a 
function $f(z) = \sum _{k=0} ^\infty  a_k z^k$.  Let $s_n(z) = \sum _ {k=0} ^ n  a_k z^k$.  The Ces\`{a}ro means $\sigma _n$  
are defined by $\sigma _n(z) = \frac{1} {n+1} \sum _{k=0} ^ n  s_k(z)$.  We note that $s_n$ can be written as the 
linear combination 

\begin{equation}
s_n = (n+1) \sigma _n - n \sigma _{n-1}.
\end{equation}

	The proof of Jentzsch's theorem (\cite{Di}, p. 352) can be modified to show that if an assumption is made that 
$\lim  \sup \sqrt [n] {| \sigma _n (z)|} \leq 1$  for some $z$ outside the circle of convergence, then a contradiction 
ensues.  Specifically, for such a $z$ and for any  $\rho  > 0$  there exists an $n(\rho )$ such that for $n > n(\rho  )$ 
we have that $\sqrt [j] {| \sigma _j(z)|} \leq (1 +\rho )$ for$ j = n-1$ and $n$.  It follows from $(2)$ that 

\begin{eqnarray}{\nonumber}
|s_n(z)| &\leq&  (n+1)|\sigma _n(z)|+n|\sigma _{n-1}(z)| \\
&\leq& (n+1)(1+\rho )^n + n(1+\rho )^{n-1} \leq 2(n+1)(1+\rho )^n.  
\end{eqnarray}

Thus, 
$\lim \sup |s_n(z)| \leq \lim \sup \sqrt [n] {2(n+1)} (1+\rho ) =1+\rho $.  The arbitrariness of $\rho$  implies that 
$\lim \sup |s_n(z)| \leq  1$, which contradicts a conclusion in the proof of Jentzsch's Theorem.  Continuing as in the proof of 
Jentzsch's Theorem, it follows that the circle of convergence is also a subset of the limit set of the zeros of 
$\{\sigma _n\}$. 



	We now consider the Ces\`{a}ro approximants of outer functions which are the derivatives of a convex functions.  
We introduce the following notation (see \cite{Du2}).  Let

\medskip

\begin{tabular}{l}
	$A =  \{ f : f(z) = a_0 + a_1 z + a_2 z^2 + . . .$,  $f$ is analytic on $\Bbb D\}$, \\
	$S = \{ f \in  A :  f(0) = 0, f (0) = 1$,  $f$ is univalent on $\Bbb D\}$,  \\
	$K = \{ f  \in S : f$ is close-to-convex $\}$,  \\
	$S^* = \{ f  \in S : f$ is starlike w.r.t. $0 \}$,  \\
            $C = \{ f \in  S : f$ is convex$\}$.  \\
\end{tabular}

\medskip

	We note Kaplan's relationship that $f  \in K$ if and only if there exists a $g \in  S^*$ such that $zf (z) = g(z)p(z)$ for 
some $p \in  A$
 such that $p(0) = 1$ and $p(\Bbb D)$ lies in a half plane $H$ with $0 \in   \partial H$.  Also, we note that close-to-convexity and convexity are 
geometric conditions on a domain which are independent of the normalization of an associated function mapping $\Bbb D$ 
onto the domain.

	For $f,g \in  A$ with $f(z) = a_0 + a_1 z + a_2 z^2 + . . .$ and $g(z) = b_0 + b_1 z + b_2 z^2 + . . .$, define 
$f*g$ by $f*g(z) = a_0 b_0 + a_1 b_1 z + a_2 b_2 z^2 + . . .$ , i.e.,  $f*g$ is the (Hadamard) convolution of $f$ and $g$.   
We note that we will employ, in the arguments we give, a common abuse of notation which
 interchanges the function $f*g$ with the function values $f(z)*g(z)$.  The two major Sheil-Small-Ruscheweyh
 results (see \cite{Du2}, Section 8.3) state that 

\medskip

\begin{tabular}{l}
	(A) 	if $f \in  C$ and $g \in  S^*$, then $f*g \in  S^*$ and \\
	(B) 	if $f \in  C$ and $g \in  S^*, p \in  A$ with $p(0) = 1$, then $f*gp = (f*g)p_1$  \\ where $p_1(\Bbb D)  \subset $ closed convex hull of $p(\Bbb D)$.  
\end{tabular}

\medskip

	Let $h \in  A$ with $h(z) = a_0 + a_1 z + a_2 z^2 + . . . $.  From the partial sums 
$s_n(z) = \sum _ {k=0} ^ n a_k z^k$ we can construct the Ces\`{a}ro means  $\sigma _n$ of $h$ by 

\begin{eqnarray}{\nonumber}
\sigma _n(z) &=& \frac{1}{n+1} \sum _{k=0} ^ n s_k(z) \\
&=&  h(z) * \sum _{k=0} ^ n  \frac{n-k+1} {n+1} z^k ~=~ h*g_n(z)
\end{eqnarray}

\noindent where 
$g_n(z) = \sum _{k=0} ^n  \frac  {n-k+1} {n+1} z^k$.  Note that $g_n$ is the Ces\`{a}ro mean of the identity function 
w.r.t. convolution, i.e., $z/(1-z)$.

	We have now

\begin{theorem} Let $f \in   A$  be such that $f(\Bbb D)$ is convex.   Then, the Ces\`{a}ro means  $\sigma _n$ of $f'$   are zero-free on 
$\Bbb \Bbb D$ for all $n$.  In particular, if $f(\Bbb D)$ is a bounded convex domain, then for the outer function $f'$   we have that the 
Ces\`{a}ro means $\sigma  _n$ are zero-free on $\Bbb D$ for all $n$.
\end{theorem}

	Proof:  Let $h =  f'$  .  Let $k$ be defined by $k(z) = z/(1-z)^2$ and note that $zf'(z) = k*f(z)$.  Then,

\begin{eqnarray}{\nonumber}  
\sigma _n(z) &=& h*g_n(z) = f '*g_n(z) = 
\frac{zf'(z)*zg_n(z)} {z} \\ &=&
\frac{f(z)*k(z)*zg_n(z)} {z} = 
\frac{f(z) * z(zg_n(z))'} {z}.
\end{eqnarray} 

If it can be shown that $zg_n  \in K$, then applying (A) and using Kaplan's relationship would yield a $g \in   S^*$ and a $p$ [with 
 $p(0) = 1$ and $p(\Bbb D)$ lying in a half plane $H$ with $0 \in   \partial H$]  such that  

$$\frac{f(z) * z(zg_n(z))'} {z} = 
\frac{f*gp(z)} {z} = 
\frac{(f*g(z))p_1(z)} {z}  \neq 0 $$ 

\noindent since $p_1(z) \neq  0$ [since $p_1(\Bbb D) \subset  H$] and $f*g \in  S^*$ is $0$ only for $z = 0$.  

	J. Lewis's (\cite{Le}, Lemma 3, p. 1118) result on Jacobi polynomials implies that $zg_n  \in K$.  An alternate approach 
to this can be made by noting that Egervary \cite{Eg}  showed that $G_n(z) = zg_n(z)$ is starlike w.r.t. $G_n(1)$.  Since another 
characterization of $K$ is that $f \in  K$ if and only if the complement of $f(\Bbb D),  \Bbb C \backslash f(\Bbb D)$, can be written as the union of 
non-intersecting half-rays, then the fact that $G_n(\Bbb D)$ is starlike with respect to $G_n(1)$ implies that $\Bbb C \backslash G_n(\Bbb D)$ is the union 
of half-rays emanating from $G_n(1)$.   

\section{\large Cone Condition}

	We now show that the ranges of the Ces\`{a}ro approximants satisfy a cone-like condition on the boundary.  
We will recall the following additional notation from geometric function theory.  Let

\medskip
\begin{tabular}{l}
	$P = \{ p \in  A : p(0) = 1$, Re $p(z) > 0$ for $z \in  \Bbb D \}$, \\
	$P(1/2) = \{ p \in  P :  $ Re  $p(z) > 1/2$ for $z \in  \Bbb D \}$, \\
	$S^*( \alpha  ) = \{  f \in  S^* :$ Re $zf' (z)/f(z) > \alpha$ for $z \in \Bbb D \} $.
\end{tabular}
\medskip

Recall (see \cite{Du2}) that $h  \in  C$ implies that $h(z)/z \in  P(1/2)$ and $zh' (z)/h(z)  \in  P(1/2)$.  Ruscheweyh (\cite{Ru}, p. 55) has 
generalized the principal results (A) and (B) on convolution to show 

\medskip
	(C)	if $f,g \in  S^*(1/2), p \in  A$ with $p(0) = 1$, then $f*gp = (f*g)p_1$  \\ where $p_1(\Bbb D)�\subset$  closed convex hull of $p(\Bbb D)$.
\medskip

\noindent Note (C) also implies that $f*g \in  S^*(1/2)$.

	We have then 

\begin{theorem}Let $f \in  A$ such that  $f(\Bbb D)$ is convex and  $f(z) = a_1 z + a_2 z^2 + . . . , a_1 \neq  0$.  Then, the Ces\`{a}ro 
means  $\sigma _n^{(2)}$ of $f'$  of order $2$ have their ranges contained in a cone (from $0$) with opening $2 \beta _n  \pi$, where  $\beta _n < 1$.  
\end{theorem}

	Proof:  We may assume that $a_1 = 1$, otherwise apply the argument to $f(z)/ a_1$.  

	First, note that for the function $f$  the Ces\`{a}ro means of $f'$  of order $\gamma$  can be given by (see \cite{Ru}, p. 142])  $f *g_n^{( \gamma )}$ where

$$g_n^{(\gamma)} (z) = \sum _{k=0} ^n  \frac { \left ( {\begin{array}{c} n-k+\gamma \\ n-k \end{array}} \right ) }  
{ \left ( {\begin{array}{c} n+\gamma \\ n \end{array}} \right ) } z^k$$.  

In Egervary's notation $s_{n-1}^{(\gamma  )} = zg_n^{(\gamma )}$.  Egervary proved for each $n \ge  1$

\medskip
\begin{tabular}{l}
	(a) $s_n^{(1)}(\Bbb D)$ is starlike w.r.t. $s_n^{(1)}(1)$, \\
	(b) $s_n^{(2)} \in   S^*(1/2)$, \\
	(c) $s_n^{(3)} \in  C$.
\end{tabular}
\medskip

	We will use (b) as follows to yield the cone condition on the ranges ${\nobreak f *g_n^{(\gamma )}(\Bbb D)}$.  Suppose $f \in  C \subset  S^*(1/2)$.  Then, 
using (b) and (C) we have 

\begin{eqnarray}{\nonumber}
f' *g_n^{(2)} &=& 
\frac{zf'(z) * z g_n^{(2)}(z)} { z} = 
\frac{f(z) * z(zg_n^{(2)}(z))'} { z}\\
&=& 
\frac {f(z) * \frac{zg_n^{(2)}z(zg_n^{(2)}(z))'}{zg_n^{(2)}}} {z} = 
\frac{f(z)*zg_n^{(2)}(z)} {z} p_1(z) 
\end{eqnarray} 

\noindent where $p_1 \in  P(1/2)$ and $f(z)*zg_n^{(2)}(z) \in  S^*(1/2)$.  By a result of Brickman, et al., \cite{BHMW} we have 
$\frac{f(z) * zg_n^{(2)}(z)} {z} \in P( 1/2 )$.  Since $(f(z)*zg_n^{(2)})/z$ is a polynomial it is bounded, hence there exists  $\beta _1 < 1$ 
such that 

$$ \left | \arg  \frac{f(z) * zg_n^{(2)}(z)} {z} \right | < \frac {\beta_1 \pi} {2}.$$ 

\noindent Hence, we have 
	
\begin{eqnarray}{\nonumber} 
\left | \arg  f'(z) * g_n^{(2)}(z)  \right | &=& 
\left | \arg \frac{f(z)*zg_n^{(2)}(z)} {z} . p_1(z) \right | \\
& \leq& 
\left | \arg \frac{f(z) * zg_n^{(2)}(z)} {z} \right | + 
| \arg p_1(z) | \\
&\leq&   
\beta _1  \frac{\pi} {2} + \frac{\pi} {2} = \beta \pi  .
\end{eqnarray}

\noindent Since a subordination argument can be made to show that   
$\frac{f(z) * zg_n^{(2)}(z)} {z} \prec g_n^{(3)}(z)$, then  $\beta  _1$ can be chosen independently of $f$.   

	We note that the above conclusion is valid for $f \in  S^*(1/2)$ and that it is not generally extendable to $f \in  S^*$, 
since $k'(z)*g_n^{(2)}$ does not satisfy a cone condition and $k�\in  S^*$.


\begin{thebibliography}{9}
\bibitem[BHMW]{BHMW}L.  Brickman, D.J. Hallenbeck, T.  MacGregor and D.R. Wilken, ``Convex hulls and extreme points 
of families of starlike and convex mappings," Trans.  Amer.  Math Soc.  185 (1973), 413-428.
\bibitem[Cu]{Cu}W.J. Cunningham, {\it Introduction to Nonlinear Analysis}, McGraw-Hill, New York, 1958.
\bibitem[Di]{Di}P.  Dienes, {\it The Taylor Series}, Dover, New York, 1957.
\bibitem[Du1]{Du1}P.  Duren, {\it Theory of $H^p$ Spaces}, Academic Press, New York, 1970.
\bibitem[Du2]{Du2}	P.  Duren, {\it Univalent Functions}, Springer-Verlag, New York, 1983.
\bibitem[Eg]{Eg}E. Egervary, ``Abbildungseigenschaften der arithmetischen Mittle der geometrischen Reihen," Math.  Z.  42 (1937), 221-230.
\bibitem[Le]{Le}J.L. Lewis, ``Applications of a convolution theorem to Jacobi polynomials," SIAM J.  Math.  Anal.  10 (1979), 1110-1120.
\bibitem[Po]{Po}C.  Pommerenke, {\it Boundary Behavior of Conformal Maps}, Springer-Verlag, New York, 1972.
\bibitem[Ru]{Ru}S.  Ruscheweyh, {\it Convolutions in Geometric Function Theory}, Sem.  Math.  Sup.,  University of Monteral Press, 83 (1982).
\end{thebibliography}

\end{document}
